% Generated draft: requires mathematical and source review.
\documentclass[11pt]{article}
\usepackage[margin=25mm]{geometry}
\usepackage{fontspec}
\setmainfont{Latin Modern Roman}
\usepackage{amsmath,amssymb,mathtools,unicode-math}
\setmathfont{Latin Modern Math}
\usepackage{enumitem,needspace,xurl,hyperref}
\setlength{\parindent}{0pt}
\setlength{\parskip}{5pt}
\setlength{\emergencystretch}{3em}
\newcommand{\N}{\mathbb{N}}
\newcommand{\Z}{\mathbb{Z}}
\newcommand{\Q}{\mathbb{Q}}
\newcommand{\R}{\mathbb{R}}
\newcommand{\C}{\mathbb{C}}
\newcommand{\F}{\mathbb{F}}
\newcommand{\A}{\mathbb{A}}
\newcommand{\PP}{\mathbb{P}}
\newcommand{\E}{\mathbb{E}}
\newcommand{\eps}{\varepsilon}
\newcommand{\dd}{\,\mathrm{d}}
\begin{document}
\section*{UnsolvedMath ID: 30003177}
\textbf{Source problem code:} catalogue entry
\subsection{Definitions and mathematical statement}
The supplied record describes a family of simple closed curves \(\Gamma_n\subset\partial(B\times I)\), each consisting of two nonparallel diameters joined by two helical arcs winding approximately \(n\) times. The source does not specify the exact family, the meaning of approximately, or the range of \(n\).

Under the usual Euclidean interpretation, \(B\) is a closed round disk in \(\R^2\), \(I=[a,b]\) with \(a<b\), and \(C=B\times I\subset\R^3\) carries the Euclidean metric. The two diameters would lie in the two end faces, with the helical arcs on the lateral surface. This is an interpretation of the description, not an explicit specification supplied by the record; in particular, the angle between the diameters, the endpoint pairing, and the helical pitches are not given.

For a specified simple closed curve \(\Gamma\subset\R^3\), define an embedded minimal disk spanning \(\Gamma\) to be a subset \(D\subset\R^3\) that is the image of a topological embedding of the closed unit disk, whose interior is a smooth embedded surface with mean curvature zero, and whose boundary is \(\Gamma\). Distinct disks are counted as distinct image subsets, rather than as different parametrizations. Write
\[
\mathcal D(\Gamma)=\{D\subset\R^3:D\text{ is an embedded minimal disk with }\partial D=\Gamma\},
\qquad N(\Gamma)=\#\mathcal D(\Gamma).
\]
This definition does not presume finiteness or smoothness across the junctions of the boundary arcs. The record does not specify boundary regularity or explicitly state whether disks must lie in \(C\).

\textbf{Problem as supplied.} Determine the embedded minimal disks spanning the described curves \(\Gamma_n\), and determine their number. In particular, does that number grow at least like \(2^n\)?

A possible precise interpretation of the growth question, after specifying a family indexed by positive integers and the admissible disks, would be
\[
\exists c>0\;\exists n_0\in\N\;\forall n\in\N,\qquad
n\ge n_0\ \Longrightarrow\ N(\Gamma_n)\ge c\,2^n.
\]
The source does not specify this interpretation, a literal bound \(N(\Gamma_n)\ge 2^n\), or another asymptotic convention. Nor does it specify whether a bound is sought for one chosen family or uniformly over all curves satisfying the informal description. Consequently, the displayed asymptotic formulation is a candidate interpretation, not an unqualified restatement.

The supplied dated assessment says that existence, parity information, and exponentially many embedded disks are known for these helicoidal boundaries, while exact classification and counting remain unresolved. It gives no quantitative theorem or hypotheses that reconcile that assessment with the exponential-growth question. The supplied open label therefore does not independently establish which precise assertion remains open.
\subsection{Short English statement}
For a cylinder boundary made from two nonparallel diameters and two winding helical arcs, describe and count the embedded minimal disks spanning it, especially whether their number grows at least like \(2^n\). The exact curve family, winding convention, admissible disks, and growth quantifiers are not specified. The supplied open status and claims about known exponential multiplicity require independent review.
\subsection{Sources}
\begin{itemize}
\item Local UnsolvedMath registry: supplied complete record with permanent ID \(30003177\), source problem code \texttt{catalogue entry}, title \emph{Exponential Multiplicity of Minimal Disks Spanning Helical Curves}, in the Oberwolfach Reports Open Problems collection. No individual proposer is supplied; the proposed year is recorded as 2016.
\item Supplied source citation: \emph{Calculus of Variations (2017)}, Oberwolfach Reports, pp.\ 1943--2008, DOI \texttt{10.4171/owr/2016/34}. \url{https://doi.org/10.4171/owr/2016/34}. The background describes Oberwolfach Reports as records of mathematical workshops and their open-problem discussions. The citation title contains 2017, whereas the DOI and proposed-year metadata contain 2016; this discrepancy is retained without resolution.
\item Supplied literature source: Hoffman--White, \emph{Number of minimal surfaces with a given boundary}. \url{https://arxiv.org/abs/0807.0933}.
\item Supplied literature source: \emph{Oberwolfach Report on calculus of variations}. \url{https://doi.org/10.4171/owr/2016/34}.
\item Literature assessment checked at \texttt{2026-08-21}: supplied dated curation, not independently verified. It labels the problem open and reports known existence, parity information, and exponentially many embedded disks, while describing exact classification and counting as unresolved. The background explicitly requests independent verification of time-sensitive status and cited claims.
\item The registry labels the statement \texttt{corrected\_verified} and says mathematical content was preserved while notation, typography, or omitted definitions were repaired for self-contained publication. These are supplied editorial assertions; the cited sources have not been independently consulted for this draft.
\end{itemize}
\end{document}
